\documentclass[11pt,a4paper]{article}
\usepackage[utf8]{inputenc}
\usepackage[english]{babel}
\usepackage{amsmath}
\usepackage{amsfonts}
\usepackage{booktabs}
\usepackage{hyperref}
\usepackage{listings}
\usepackage{xcolor}
\usepackage{geometry}

\geometry{margin=1in}

\definecolor{codegreen}{rgb}{0,0.6,0}
\definecolor{codegray}{rgb}{0.5,0.5,0.5}
\definecolor{backcolour}{rgb}{0.95,0.95,0.92}

\lstset{
    backgroundcolor=\color{backcolour},
    commentstyle=\color{codegreen},
    keywordstyle=\color{magenta},
    numberstyle=\tiny\color{codegray},
    stringstyle=\color{purple},
    basicstyle=\ttfamily\small,
    breakatwhitespace=false,
    breaklines=true,
    captionpos=b,
    keepspaces=true,
    numbers=left,
    numbersep=5pt,
    showspaces=false,
    showstringspaces=false,
    showtabs=false,
    tabsize=2,
    frame=single
}

\title{Code Optimization Report: Volumetric Solidification \\ \large Architectural Enhancements for Haskell-Vulkan \texttt{Clouds.hs}}
\author{Advanced Graphics Engine Group}
\date{\today}

\begin{document}

\maketitle

\section{Diagnostic: The Multi-Layer Slicing Phenomenon}
When volumetric media resembles discrete parallel slices rather than a contiguous solid structure, it indicates that the raymarch step increment matches or beats the spatial frequencies of the 3D texture lookup. 

\begin{enumerate}
    \item \textbf{Deterministic Phase Alignment:} If view rays iterate over fixed spatial deltas ($\Delta t$), every neighbor screen pixel samples identical offsets relative to the voxel grid boundaries. This shifts the perceptual shape into a series of transparent topological shell steps.
    \item \textbf{Missing Spatial Interleaving:} Without randomizing the origin or evaluation points along individual projection arcs, the underlying voxel texture bounds project flat lines over perspective depth changes.
\end{enumerate}

To turn these shells into solid shapes, we inject \textbf{Blue Noise Interleaved Jittering} into the entry point calculation and leverage \textbf{Temporal Reprojection Matrices} to smooth out the grain across frames.

---

\section{Mathematical Matrix Operations for Temporal Reprojection}
To achieve high-fidelity solidification without massive sample counts, the shader scales to quarter-resolution and reprojects history. This relies on tracking individual positions backwards through time via custom matrix operations.

\subsection{Matrix Reprojection Pipeline}
Let $P_{\text{curr}}$ be the reconstructed 3D World Space position where the ray terminates or samples a cloud boundary during the current frame. To retrieve the matching pixel coordinate from the previous frame ($UV_{\text{prev}}$):

\begin{equation}
    X_{\text{clip\_prev}} = \mathbf{M_{\text{PrevViewProj}}} \cdot \begin{bmatrix} P_{\text{curr}.x} - W_{\text{offset}.x} \\ P_{\text{curr}.y} \\ P_{\text{curr}.z} - W_{\text{offset}.z} \\ 1.0 \end{bmatrix}
\end{equation}

Where $\mathbf{M_{\text{PrevViewProj}}}$ is the non-reversed view-projection transform tracking the position of the camera in the previous frame, and $W_{\text{offset}}$ accounts for the global translation displacement induced by the wind velocity vector accumulated across the delta time interval $\Delta t$.

\subsection{Perspective Projection Space Normalization}
The tracking space must transition from Homogeneous 4D Clip Coordinates to Normalized Device Coordinates (NDC) via the perspective divide:

\begin{equation}
    X_{\text{ndc\_prev}} = \frac{X_{\text{clip\_prev}.xyz}}{X_{\text{clip\_prev}.w}}
\end{equation}

Finally, map the range $[-1.0, 1.0]^2$ into standard screen coordinate space $[0.0, 1.0]^2$ to fetch historical data inputs:

\begin{equation}
    UV_{\text{prev}} = X_{\text{ndc\_prev}.xy} \cdot \begin{bmatrix} 0.5 \\ 0.5 \end{bmatrix} + \begin{bmatrix} 0.5 \\ 0.5 \end{bmatrix}
\end{equation}

---

\section{Haskell FIR Code Implementation Example}
Below is the optimized source block rewritten for the Haskell \texttt{FIR} Vulkan pipeline DSL. It includes **Spherical Earth Curvature**, **Dynamic Jittering offsets**, and **Nested Raymarching structures**.

\begin{lstlisting}[language=Haskell, caption={Haskell FIR DSL Implementation for Curved Volumetric Shells}]
{-# LANGUAGE DataKinds #-}
{-# LANGUAGE TypeOperators #-}

module Graphics.Haskan.Vulkan.Shaders.Deferred.Clouds where

-- Conceptual representation of the dynamic raymarching pipeline stage
type CloudShaderPipeline = '[
    "world_pos"    ':-> Input  (Vec 3 Float),
    "blue_noise"   ':-> Texture2D (Vec 4 Float),
    "cloud_noise"  ':-> Texture3D (Vec 4 Float),
    "out_color"    ':-> Output (Vec 4 Float)
  ]

-- Realization of curved earth surface offsets and dynamic loop execution
traceCurvedClouds :: Program CloudShaderPipeline ()
traceCurvedClouds = do
    camPos   <- use @"camera_pos"
    viewDir  <- use @"view_direction"
    uvCoords <- use @"screen_uv"
    
    -- Fetch spatial blue noise to break slice layers
    noiseSample <- sample @"blue_noise" uvCoords
    let jitter = noiseSample .& R -- Extract red channel seed scalar

    -- Earth planetary scale dimensions
    let earthRadius    = 6371000.0
    let cloudBottom    = 2500.0
    let cloudThickness = 800.0
    
    -- Spherical atmospheric intersection logic
    let camHeightCurved = camPos .& Y + earthRadius
    let tNear = computeSphericalShellIntersection camHeightCurved viewDir cloudBottom
    
    -- Inject stochastic jitter into starting path
    let stepSize = 16.0
    let startPos = camPos + viewDir * (tNear + (jitter * stepSize))
    
    -- Dynamic Loop variable configuration
    accumOpacity <- lit 0.0
    accumLight   <- lit (Vec3 0.0 0.0 0.0)
    rayDistance  <- lit 0.0
    
    loop do
        -- Terminate if ray leaves bounds or opacity saturates
        isSaturated  <- accumOpacity >=. 0.99
        isPastVolume <- rayDistance  >=. cloudThickness
        
        breakIf (isSaturated ||. isPastVolume)
        
        let currPos = startPos + viewDir * rayDistance
        
        -- Apply Curvature Matrix shift along horizontal trajectory planes
        let distHorizSq = (currPos.&X - camPos.&X)^2 + (currPos.&Z - camPos.&Z)^2
        let curvedY     = currPos.&Y - (distHorizSq / (2.0 * earthRadius))
        
        -- Sample Density Mapping with height mask filters
        let heightNorm = (curvedY - cloudBottom) / cloudThickness
        density <- sampleDensityField currPos heightNorm
        
        when (density >. 0.0) $ do
            -- In-scatter illumination calculations (Beer-Powder integration)
            transmittance <- calculateBeerPowderLight currPos density
            accumLight   .= accumLight + (transmittance * density * stepSize)
            accumOpacity .= accumOpacity + ((1.0 - accumOpacity) * density * stepSize)
            
        rayDistance .= rayDistance + stepSize
        
    write @"out_color" (Vec4 (accumLight.&X) (accumLight.&Y) (accumLight.&Z) accumOpacity)
\end{lstlisting}

---

\section{Summary of Parameter Configuration Shifts}

\begin{table}[h!]
\centering
\small
\begin{tabular}{lll}
\toprule
\textbf{Pipeline Element} & \textbf{Layered/Slicing Setup} & \textbf{Solidified Volume Setup} \\ \midrule
\textbf{Ray Start Offset} & Absolute \texttt{tNear} Step Grid & Blue Noise Stochastic Ray Jitter \\
\textbf{Atmosphere Plane} & Flat Box Bounds Matrix & Curved Earth Coordinates Displacement \\
\textbf{History Filter} & Raw Texel Blending & Velocity and Wind Corrected Matrix Pull \\
\textbf{Light Absorption} & Beer's Law (Linear Linear) & Beer-Powder Attenuation Curve \\ \bottomrule
\end{tabular}
\caption{Architectural changes implemented to enforce volume cohesion.}
\end{table}

\end{document}

